Die Gewichtskraft ist das Produkt aus Masse und Fallbeschleunigung. Auf der Erde ist
\begin{align}
 F_\text{G,E} &= \FGEF = \mK\cdot\gE = \FGE \approx \result{\FGEP{0}{2}}
\end{align}

Die Masse ist eine Eigenschaft des Körpers und hängt nicht vom Ort ab. Auf dem Mond und auf dem Mars ist sie gleich gross wie auf der Erde:
\begin{align}
 m &= \result{\mKrP{0}{2}}
\end{align}

Mit den Fallbeschleunigungen $g_\text{Mo} = \gMoO$ (Mond, FoTaBe S.~210) und $g_\text{Ma} = \gMaO$ (Mars, FoTaBe S.~212) folgt
\begin{align}
 F_\text{G,Mo} &= \FGMoF = \mK\cdot\gMo = \FGMo \approx \result{\FGMoP{0}{2}} \\
 F_\text{G,Ma} &= \FGMaF = \mK\cdot\gMa = \FGMa \approx \result{\FGMaP{0}{2}}
\end{align}
Auf dem Mond wiegt man also nur etwa einen Sechstel, auf dem Mars gut einen Drittel so viel wie auf der Erde.
